% Explanation certificate for heart-s5-c3 (genome 96fb45fb-972d-4aaa-b6d5-01fc63fbc4e6).
% Generated by explaneat export (schema 1.4.0) from the JSON record alone;
% content sha256 644db15ecc0191f8ed73b9a23757dccfcf5f1809df487f9171d22800751515e6.
% Drafter recorded in the JSON: ai_agent (claude). Do not edit by hand: re-render with
%   uv run python -m explaneat export render --from heart-s5-c3.json
% TikZ modules are input from \explaneatcertdir{<label>} (empty by default: same
% directory; index.tex sets it to the genome's subdirectory).
\providecommand{\explaneatcertdir}[1]{}
\begin{drafted}[appendix: heart-s5-c3 explanation certificate]
\section{Explanation certificate: heart-s5-c3}\label{app:heart-s5-c3:sec}
This section is the explanation certificate of the heart-s5-c3 network. It is rendered from the canonical JSON record \texttt{heart-s5-c3.json} (schema 1.4.0, SHA-256 \texttt{644db15ecc0191f8}\ldots), which also holds every evidence payload, the full graphs with weights, and the explanation record. \Cref{app:heart-s5-c3:summary} summarises the network and its coverage.

\begin{table}[htbp]
\centering\small
\begin{tabularx}{\linewidth}{@{}lX@{}}
\toprule
Dataset & Heart Disease (binary, prepared) \\
Experiment & 10seed-heart-s5, seed 5 \\
Genome & \texttt{96fb45fb-972d-4aaa-b6d5-01fc63fbc4e6} (NEAT key 9698, generation 39) \\
Explanation & \texttt{1b3f669d-9fb5-4126-b6e4-35535db7b3b1} (30 operations) \\
Train AUC / accuracy & 0.9376 / 0.8884 (n = 242) \\
Test AUC / accuracy & 0.9084 / 0.8197 (n = 61) \\
Output activation & sigmoid \\
Phenotype & 31 nodes (26 inputs, 4 hidden, 1 output), 26 enabled edges \\
Explained model state & 32 nodes, 26 enabled edges \\
Refactoring operations & split\_\allowbreak{}node 1 \\
Function preservation & max $|\Delta\,\mathrm{output}|$ = $2.22\times 10^{-16}$ on the training split \\
Annotations & 3 (2 leaf, 1 composition) \\
Structural coverage & 1 (31 of 31 candidate nodes; 7 Null) \\
Compositional coverage & 1 \\
Uncovered nodes / edges & 0 / 0 \\
Missing compositions & 0 \\
Full coverage & yes \\
Drafted by & AI-assisted draft \\
Verified by & unverified \\
Record hash (SHA-256) & \texttt{644db15ecc0191f8}\ldots \\
\bottomrule
\end{tabularx}
\caption{Explanation certificate summary: heart-s5-c3.}
\label{app:heart-s5-c3:summary}
\end{table}

\paragraph{Drafter's summary.} Mostly logistic regression: 17 direct edges (logit SD 3.22; train AUC 0.932 alone) plus an ST-depression chain and an age hinge (logit SD 0.80 together) to 0.938. The chain scores oldpeak (threshold lowered by a normal resting ECG) and adds a second, much steeper stage of the same score, so oldpeak's effect is a hinge (about 0.12 logit per SD, then about 1.10 per SD above z -0.39 with a normal ECG or +0.40 without); age enters only through a small hinge starting just below the median. Same ischaemia gate as the 10-seed heart-s5 record (key 9330) and its cluster-1 sibling, but with no resting-BP offset and a free-standing age term.

\paragraph{What resisted explanation.} Nothing resisted coverage. The chain 204 -$>$ 3487 -$>$ 290 (with 204 -$>$ 290 as well, and both 204 and 290 read out) was kept as one two-exit leaf rather than two leaves: 3487's only input is 204, so it cannot be intercepted by an identity node (add\_\allowbreak{}identity\_\allowbreak{}node must leave one input un-redirected), and splitting 204 would have duplicated the gate. 3487 is a near-copy of 204 (ReLU(0.94*204 - 0.002)), so 290 = ReLU(1.818*204 - 0.470): a second, steeper threshold on the same score. Counter-directional signs reported, not explained: a normal resting ECG (restecg:0) lowers the oldpeak threshold inside the chain (raising risk) while its direct edge is -0.188; fbs:0 +0.27 directly; thal:7 only +0.025. The age hinge barely matters (readout ablation -0.0005 AUC). closed\_\allowbreak{}form\_\allowbreak{}max\_\allowbreak{}abs\_\allowbreak{}err 0.021 comes from 3-significant-figure display rounding in the nested chain formula (numeric error 1.8e-15).

\subsection{Plain-language summary}\label{app:heart-s5-c3:plain}
\noindent{\small\textit{AI-assisted draft for review (unverified). Each sentence cites the claims (C01, C02, \ldots) that back it; every cited claim was computed on the model and the training split and passed its automatic checks. Magnitudes are changes in predicted risk for a typical patient; confidence words follow the test-split check and n. The text describes what the model does, not the world.}}\par\smallskip
\paragraph{Summary (drafted).} For a typical patient (a 56-year-old man with no chest pain symptoms and median test results) the model predicts a risk of heart disease of about 18\%; it adds up most of its inputs independently, and two parts add to that sum: a signal from ST depression on exercise and the resting ECG that acts on 105 of 242 training patients, and an age term that acts on 148.\,{\scriptsize\textcolor{NavyBlue}{[C25, C30, C34]}} These add up independently, and these lower the risk: being a woman takes a typical patient's predicted risk from about 18\% to about 2\%, typical angina to about 2\%, non-anginal chest pain to about 4\% and atypical angina to about 10\%; fasting blood sugar above 120 mg/dL lowers it to about 9\%; and a higher maximum heart rate lowers it steadily, from about 23\% at 118 to about 15\% at 175 beats per minute.\,{\scriptsize\textcolor{NavyBlue}{[C10, C09, C11, C15, C12, C14]}} Each extra vessel coloured by fluoroscopy adds about the same, from about 18\% at 0 to about 86\% at 2 vessels, flattening out near the top because the predicted risk cannot go much higher; a flat rather than upsloping ST segment raises it to about 56\%, a reversible defect on the thallium scan to about 38\% and chest pain brought on by exercise to about 35\%; higher cholesterol raises it steadily, from about 11\% at 193 mg/dL to about 33\% at 308 mg/dL, and higher resting blood pressure steadily, from about 12\% at 110 mm Hg to about 30\% at 155 mm Hg.\,{\scriptsize\textcolor{NavyBlue}{[C01, C04, C06, C08, C05, C07]}} Through this sum alone, a resting ECG other than normal raises it from about 18\% to about 24\%; the resting ECG also acts through the ST depression signal (below).\,{\scriptsize\textcolor{NavyBlue}{[C16]}} Some rarer results also count (the model has little data on this, n = 2 to 17): a thallium result coded 4 raises the risk to about 79\%, a fixed defect on the thallium scan lowers it to about 9\%, and a downsloping ST segment leaves it at about 19\%.\,{\scriptsize\textcolor{NavyBlue}{[C02, C13, C19]}} This part reads ST depression on exercise and the resting ECG, and acts on 105 of 242 training patients, 66 of whom had heart disease; for the typical patient it acts on, the predicted risk is about 76\% without it and about 91\% with it.\,{\scriptsize\textcolor{NavyBlue}{[C30]}} Through it, more ST depression adds nothing until about 0.5 mm, then raises the predicted risk steadily: for a typical patient from about 16\% at 0 mm to about 60\% at 2.8 mm.\,{\scriptsize\textcolor{NavyBlue}{[C26]}} ST depression counts more with a normal resting ECG, nothing until about 0.5 mm and then steadily, from about 16\% to about 60\%, against about 21\% to about 48\% with any other resting ECG.\,{\scriptsize\textcolor{NavyBlue}{[C28]}} Through it, a resting ECG other than normal lowers a typical patient's risk from about 18\% to about 16\%.\,{\scriptsize\textcolor{NavyBlue}{[C27]}} This part acts on 148 of 242 training patients, 81 of whom had heart disease; for the typical patient it acts on, the predicted risk is about 91\% without it and about 92\% with it.\,{\scriptsize\textcolor{NavyBlue}{[C34]}} Through it, age adds nothing until about 53 years, then raises the risk steadily, for a typical patient from about 18\% at 42 to about 19\% at 66 years.\,{\scriptsize\textcolor{NavyBlue}{[C32, C33]}}
\paragraph{The whole model.} The model adds up its inputs into one score: the vessels coloured (each extra vessel adding about the same), the ST segment, sex and chest pain type each move the predicted risk on their own, and two parts add to that sum: a signal from ST depression shaped by the resting ECG, and an age term.\,{\scriptsize\textcolor{NavyBlue}{[C25, C01, C04, C10, C09]}} \textit{Evidence.} It is made up of the inputs that feed the score directly and two parts, the ST depression signal and the age term (each explained below), and the score is their sum, which was checked on all 242 training patients.\,{\scriptsize\textcolor{NavyBlue}{[C25]}} ST depression on exercise and age do not feed the score directly, so they act only through these parts.\,{\scriptsize\textcolor{NavyBlue}{[C25]}} Changing one input at a time across the whole model, more ST depression adds nothing until about 0.5 mm, then raises a typical patient's predicted risk steadily, from about 16\% at 0 mm to about 60\% at 2.8 mm.\,{\scriptsize\textcolor{NavyBlue}{[C03]}} Across the whole model, older age adds nothing until about 53 years, then raises it steadily, from about 18\% at 42 to about 19\% at 66 years.\,{\scriptsize\textcolor{NavyBlue}{[C18]}}
\paragraph{ST depression signal.} This part raises the predicted risk with ST depression on exercise, adding nothing until about 0.5 mm and then steadily, and more so with a normal resting ECG.\,{\scriptsize\textcolor{NavyBlue}{[C26, C28]}} \textit{Evidence.} It acts on 105 of 242 training patients, 66 of whom had heart disease; for the typical patient it acts on, the predicted risk is about 76\% without it and about 91\% with it.\,{\scriptsize\textcolor{NavyBlue}{[C30]}} Through it, more ST depression adds nothing until about 0.5 mm, then raises the predicted risk steadily: for a typical patient from about 16\% at 0 mm to about 60\% at 2.8 mm.\,{\scriptsize\textcolor{NavyBlue}{[C26]}} ST depression counts more with a normal resting ECG, nothing until about 0.5 mm and then steadily, from about 16\% to about 60\%, against about 21\% to about 48\% with any other resting ECG.\,{\scriptsize\textcolor{NavyBlue}{[C28]}} Through it, a resting ECG other than normal lowers a typical patient's risk from about 18\% to about 16\%.\,{\scriptsize\textcolor{NavyBlue}{[C27]}} Its formula reproduces the part exactly on all 242 training patients.\,{\scriptsize\textcolor{NavyBlue}{[C31]}}
\paragraph{Age term.} This part adds a little to the predicted risk for older patients: nothing until about 53 years, then steadily.\,{\scriptsize\textcolor{NavyBlue}{[C32, C33]}} \textit{Evidence.} It acts on 148 of 242 training patients, 81 of whom had heart disease; for the typical patient it acts on, the predicted risk is about 91\% without it and about 92\% with it.\,{\scriptsize\textcolor{NavyBlue}{[C34]}} Through it, age adds nothing until about 53 years, then raises the risk steadily, for a typical patient from about 18\% at 42 to about 19\% at 66 years.\,{\scriptsize\textcolor{NavyBlue}{[C32, C33]}} Its formula reproduces the part exactly on all 242 training patients.\,{\scriptsize\textcolor{NavyBlue}{[C36]}}
\paragraph{Worth checking against what is known.} Worth checking: the model treats fasting blood sugar above 120 mg/dL as lowering the risk of heart disease (for a typical patient from about 18\% to about 9\%).\,{\scriptsize\textcolor{NavyBlue}{[C12]}} Worth checking: patients with typical angina get the lowest predicted risk of all chest pain types, below patients with no chest pain symptoms.\,{\scriptsize\textcolor{NavyBlue}{[C09, C11, C15]}} Worth checking: the model treats a fixed defect on the thallium scan as lowering the risk compared with a normal scan, from only 15 training patients.\,{\scriptsize\textcolor{NavyBlue}{[C13]}} Worth checking: the model gives a higher predicted risk to the 2 training patients whose thallium result is coded 4, a code the dataset documentation does not define.\,{\scriptsize\textcolor{NavyBlue}{[C02]}} Worth checking: ST depression on exercise raises the predicted risk more when the resting ECG is normal than when it is not.\,{\scriptsize\textcolor{NavyBlue}{[C28]}}
\paragraph{Claims.}
\par{\tiny\noindent
\begin{tabularx}{\linewidth}{@{}l>{\raggedright\arraybackslash}X>{\raggedright\arraybackslash}p{0.11\linewidth}>{\raggedright\arraybackslash}p{0.16\linewidth}>{\raggedright\arraybackslash}p{0.05\linewidth}>{\raggedright\arraybackslash}p{0.15\linewidth}l>{\raggedright\arraybackslash}p{0.08\linewidth}@{}}
\toprule
Claim & What is checked & Risk, typical patient & Estimate (log-odds; OR); range across patients & n & Test check & Verified & Words \\
\midrule
C01 & direction: number of vessels coloured by fluoroscopy: 0 vessels to 2 vessels (whole model) \newline\textit{shape: steady (equal steps); 100\% of patients same; verified; risk saturates near 100\%} & 18\% to 86\% & log-odds 3.33, OR 28.0; range across patients 3.33 to 3.33 & 242 & train; test unclear (AUC gain 0.029 [-0.039, 0.100], n=61) & \textcolor{OliveGreen}{yes} & on the training data only \\
C02 & direction: thallium stress test result: a normal thallium scan to thallium stress test result coded 4 (meaning not recorded) (whole model) \newline\textit{shape: categorical; highest to lowest: thallium stress test result coded 4 (meaning not recorded), a reversible defect on the thallium scan, a normal thallium scan, a fixed defect on the thallium scan} & 18\% to 79\% & log-odds 2.85, OR 17.4; range across patients 2.85 to 2.85 & 135, 2 (small n) & train; test unclear (AUC gain 0.020 [-0.012, 0.058], n=61) & \textcolor{OliveGreen}{yes} & the model has little data on this; on the training data only \\
C03 & direction: ST depression on exercise: 0.0 mm to 2.8 mm (whole model) \newline\textit{shape: hinge at 0.5 mm; 100\% of patients same; verified} & 16\% to 60\% & log-odds 1.65, OR 5.20; range across patients 1.22 to 2.08 & 242 & train; test unclear (AUC gain 0.023 [-0.001, 0.055], n=61) & \textcolor{OliveGreen}{yes} & on the training data only \\
C04 & direction: slope of the ST segment at peak exercise: an upsloping ST segment to a flat ST segment (whole model) \newline\textit{shape: categorical; highest to lowest: a flat ST segment, a downsloping ST segment, an upsloping ST segment} & 18\% to 56\% & log-odds 1.75, OR 5.75; range across patients 1.75 to 1.75 & 115, 110 & train; test disagrees (AUC gain -0.024 [-0.071, 0.009], n=61) & \textcolor{OliveGreen}{yes} & on the training data only \\
C05 & direction: serum cholesterol: 193 mg/dL (5.0 mmol/L) to 308 mg/dL (8.0 mmol/L) (whole model) \newline\textit{shape: steady; 100\% of patients same; verified} & 11\% to 33\% & log-odds 1.32, OR 3.73; range across patients 1.32 to 1.32 & 242 & train; test disagrees (AUC gain -0.019 [-0.058, 0.008], n=61) & \textcolor{OliveGreen}{yes} & on the training data only \\
C06 & direction: thallium stress test result: a normal thallium scan to a reversible defect on the thallium scan (whole model) \newline\textit{shape: categorical; highest to lowest: thallium stress test result coded 4 (meaning not recorded), a reversible defect on the thallium scan, a normal thallium scan, a fixed defect on the thallium scan} & 18\% to 38\% & log-odds 1.00, OR 2.72; range across patients 1.00 to 1.00 & 135, 90 & train; test unclear (AUC gain 0.020 [-0.012, 0.058], n=61) & \textcolor{OliveGreen}{yes} & on the training data only \\
C07 & direction: resting blood pressure: 110 mm Hg to 155 mm Hg (whole model) \newline\textit{shape: steady; 100\% of patients same; verified} & 12\% to 30\% & log-odds 1.19, OR 3.30; range across patients 1.19 to 1.19 & 242 & train; test unclear (AUC gain 0.008 [-0.018, 0.037], n=61) & \textcolor{OliveGreen}{yes} & on the training data only \\
C08 & direction: exercise-induced angina: no chest pain on exercise to chest pain brought on by exercise (whole model) \newline\textit{shape: categorical; highest to lowest: chest pain brought on by exercise, no chest pain on exercise} & 18\% to 35\% & log-odds 0.89, OR 2.44; range across patients 0.89 to 0.89 & 160, 82 & train; test agrees (AUC gain 0.022 [0.003, 0.049], n=61) & \textcolor{OliveGreen}{yes} & clearly \\
C09 & direction: chest pain type: no chest pain (asymptomatic) to typical angina (whole model) \newline\textit{shape: categorical; highest to lowest: no chest pain (asymptomatic), atypical angina, non-anginal chest pain, typical angina} & 18\% to 2\% & log-odds -2.43, OR 0.09; range across patients -2.43 to -2.43 & 108, 22 & train; test agrees (AUC gain 0.033 [0.008, 0.066], n=61) & \textcolor{OliveGreen}{yes} & clearly \\
C10 & direction: sex: male to female (whole model) \newline\textit{shape: categorical; highest to lowest: male, female} & 18\% to 2\% & log-odds -2.19, OR 0.11; range across patients -2.19 to -2.19 & 159, 83 & train; test unclear (AUC gain 0.025 [-0.007, 0.066], n=61) & \textcolor{OliveGreen}{yes} & on the training data only \\
C11 & direction: chest pain type: no chest pain (asymptomatic) to non-anginal chest pain (whole model) \newline\textit{shape: categorical; highest to lowest: no chest pain (asymptomatic), atypical angina, non-anginal chest pain, typical angina} & 18\% to 4\% & log-odds -1.69, OR 0.18; range across patients -1.69 to -1.69 & 108, 70 & train; test agrees (AUC gain 0.033 [0.008, 0.066], n=61) & \textcolor{OliveGreen}{yes} & clearly \\
C12 & direction: fasting blood sugar: fasting blood sugar of 120 mg/dL or less to fasting blood sugar above 120 mg/dL (whole model) \newline\textit{shape: categorical; highest to lowest: fasting blood sugar of 120 mg/dL or less, fasting blood sugar above 120 mg/dL} & 18\% to 9\% & log-odds -0.80, OR 0.45; range across patients -0.80 to -0.80 & 209, 33 & train; test unclear (AUC gain 0.000 [-0.013, 0.013], n=61) & \textcolor{OliveGreen}{yes} & on the training data only \\
\bottomrule
\end{tabularx}}
\par{\tiny\noindent
\begin{tabularx}{\linewidth}{@{}l>{\raggedright\arraybackslash}X>{\raggedright\arraybackslash}p{0.11\linewidth}>{\raggedright\arraybackslash}p{0.16\linewidth}>{\raggedright\arraybackslash}p{0.05\linewidth}>{\raggedright\arraybackslash}p{0.15\linewidth}l>{\raggedright\arraybackslash}p{0.08\linewidth}@{}}
\toprule
Claim & What is checked & Risk, typical patient & Estimate (log-odds; OR); range across patients & n & Test check & Verified & Words \\
\midrule
C13 & direction: thallium stress test result: a normal thallium scan to a fixed defect on the thallium scan (whole model) \newline\textit{shape: categorical; highest to lowest: thallium stress test result coded 4 (meaning not recorded), a reversible defect on the thallium scan, a normal thallium scan, a fixed defect on the thallium scan} & 18\% to 9\% & log-odds -0.80, OR 0.45; range across patients -0.80 to -0.80 & 135, 15 (small n) & train; test unclear (AUC gain 0.020 [-0.012, 0.058], n=61) & \textcolor{OliveGreen}{yes} & the model has little data on this; on the training data only \\
C14 & direction: maximum heart rate: 118 beats per minute to 175 beats per minute (whole model) \newline\textit{shape: steady; 100\% of patients same; verified} & 23\% to 15\% & log-odds -0.54, OR 0.58; range across patients -0.54 to -0.54 & 242 & train; test agrees (AUC gain 0.014 [0.002, 0.035], n=61) & \textcolor{OliveGreen}{yes} & clearly \\
C15 & direction: chest pain type: no chest pain (asymptomatic) to atypical angina (whole model) \newline\textit{shape: categorical; highest to lowest: no chest pain (asymptomatic), atypical angina, non-anginal chest pain, typical angina} & 18\% to 10\% & log-odds -0.71, OR 0.49; range across patients -0.71 to -0.71 & 108, 42 & train; test agrees (AUC gain 0.033 [0.008, 0.066], n=61) & \textcolor{OliveGreen}{yes} & clearly \\
C16 & direction: resting ECG: a normal resting ECG to resting ECG other than a normal resting ECG (direct contribution only, not through any part) \newline\textit{shape: categorical; highest to lowest: resting ECG other than a normal resting ECG, a normal resting ECG} & 18\% to 24\% & log-odds 0.38, OR 1.46; range across patients 0.38 to 0.38 & 121, 121 & train; test unclear (AUC gain 0.005 [-0.005, 0.024], n=61) & \textcolor{OliveGreen}{yes} & on the training data only \\
C17 & direction: resting ECG: a normal resting ECG to resting ECG other than a normal resting ECG (whole model) \newline\textit{shape: categorical; highest to lowest: resting ECG other than a normal resting ECG, a normal resting ECG} & 18\% to 21\% & log-odds 0.02, OR 1.02; range across patients -0.49 to 0.38 & 121, 121 & train; test disagrees (AUC gain 0.003 [-0.013, 0.021], n=61) & \textcolor{BrickRed}{no} & slightly; on the training data only \\
C18 & direction: age: 42 years to 66 years (whole model) \newline\textit{shape: hinge at 53 years; 100\% of patients same; verified; matches the fold at 53 years} & 18\% to 19\% & log-odds 0.10, OR 1.11; range across patients 0.10 to 0.10 & 242 & train; test disagrees (AUC gain -0.003 [-0.017, 0.004], n=61) & \textcolor{OliveGreen}{yes} & slightly; on the training data only \\
C19 & direction: slope of the ST segment at peak exercise: an upsloping ST segment to a downsloping ST segment (whole model) \newline\textit{shape: categorical; highest to lowest: a flat ST segment, a downsloping ST segment, an upsloping ST segment} & 18\% to 19\% & log-odds 0.05, OR 1.05; range across patients 0.05 to 0.05 & 115, 17 (small n) & train; test disagrees (AUC gain -0.024 [-0.071, 0.009], n=61) & \textcolor{OliveGreen}{yes} & slightly; the model has little data on this; on the training data only \\
C20 & magnitude: number of vessels coloured by fluoroscopy: 0 vessels to 2 vessels, typical patient \newline\textit{shape: steady (equal steps); 100\% of patients same; verified; risk saturates near 100\%} & 18\% to 86\% & log-odds 3.33, OR 28.0 & 242 & train; test unclear (AUC gain 0.029 [-0.039, 0.100], n=61) & \textcolor{OliveGreen}{yes} & on the training data only \\
C21 & magnitude: ST depression on exercise: 0.0 mm to 2.8 mm, typical patient \newline\textit{shape: hinge at 0.5 mm; 100\% of patients same; verified} & 16\% to 60\% & log-odds 2.08, OR 8.01 & 242 & train; test unclear (AUC gain 0.023 [-0.001, 0.055], n=61) & \textcolor{OliveGreen}{yes} & on the training data only \\
C22 & magnitude: slope of the ST segment at peak exercise: an upsloping ST segment to a flat ST segment, typical patient \newline\textit{shape: categorical; highest to lowest: a flat ST segment, a downsloping ST segment, an upsloping ST segment} & 18\% to 56\% & log-odds 1.75, OR 5.75 & 115, 110 & train; test disagrees (AUC gain -0.024 [-0.071, 0.009], n=61) & \textcolor{OliveGreen}{yes} & on the training data only \\
C23 & magnitude: serum cholesterol: 193 mg/dL (5.0 mmol/L) to 308 mg/dL (8.0 mmol/L), typical patient \newline\textit{shape: steady; 100\% of patients same; verified} & 11\% to 33\% & log-odds 1.32, OR 3.73 & 242 & train; test disagrees (AUC gain -0.019 [-0.058, 0.008], n=61) & \textcolor{OliveGreen}{yes} & on the training data only \\
C24 & magnitude: thallium stress test result: a normal thallium scan to a reversible defect on the thallium scan, typical patient \newline\textit{shape: categorical; highest to lowest: thallium stress test result coded 4 (meaning not recorded), a reversible defect on the thallium scan, a normal thallium scan, a fixed defect on the thallium scan} & 18\% to 38\% & log-odds 1.00, OR 2.72 & 135, 90 & train; test unclear (AUC gain 0.020 [-0.012, 0.058], n=61) & \textcolor{OliveGreen}{yes} & on the training data only \\
\bottomrule
\end{tabularx}}
\par{\tiny\noindent
\begin{tabularx}{\linewidth}{@{}l>{\raggedright\arraybackslash}X>{\raggedright\arraybackslash}p{0.11\linewidth}>{\raggedright\arraybackslash}p{0.16\linewidth}>{\raggedright\arraybackslash}p{0.05\linewidth}>{\raggedright\arraybackslash}p{0.15\linewidth}l>{\raggedright\arraybackslash}p{0.08\linewidth}@{}}
\toprule
Claim & What is checked & Risk, typical patient & Estimate (log-odds; OR); range across patients & n & Test check & Verified & Words \\
\midrule
C25 & structure: made of ischaemia\_\allowbreak{}chain, age\_\allowbreak{}hinge and inputs sex, cp, trestbps, chol, fbs, restecg, thalach, exang, slope, ca, thal; added &  & log-odds n/a, OR n/a & 242 & train & \textcolor{OliveGreen}{yes} &  \\
C26 & direction: ST depression on exercise: 0.0 mm to 2.8 mm (through this part) \newline\textit{shape: hinge at 0.5 mm; 100\% of patients same; verified} & 16\% to 60\% & log-odds 1.65, OR 5.20; range across patients 1.22 to 2.08 & 242 & train; test unclear (AUC gain 0.023 [-0.001, 0.055], n=61) & \textcolor{OliveGreen}{yes} & on the training data only \\
C27 & direction: resting ECG: a normal resting ECG to resting ECG other than a normal resting ECG (through this part) \newline\textit{shape: categorical; highest to lowest: a normal resting ECG, resting ECG other than a normal resting ECG} & 18\% to 16\% & log-odds -0.36, OR 0.70; range across patients -0.86 to 0.00 & 121, 121 & train; test unclear (AUC gain 0.006 [-0.009, 0.028], n=61) & \textcolor{OliveGreen}{yes} & slightly; on the training data only \\
C28 & interaction: effect of ST depression on exercise (0.0 mm to 2.8 mm) by resting ECG: resting ECG other than a normal resting ECG minus a normal resting ECG; stronger in a normal resting ECG \newline\textit{shape: hinge at 0.5 mm; 100\% of patients same; verified} & resting ECG other than a normal resting ECG: 21\% to 48\%; a normal resting ECG: 16\% to 60\% & log-odds difference -0.86, OR 0.42 & 121, 121 & train; test agrees (AUC gain 0.006 [-0.009, 0.028], n=61); group difference on test -0.86 [-0.86, -0.86] & \textcolor{OliveGreen}{yes} & clearly \\
C29 & threshold: ST depression on exercise: fold at about 0.9 mm (starts) \newline\textit{shape: hinge at 0.5 mm; 100\% of patients same; verified} & 16\% to 74\% (typical patient with a fold) & fold 0.92 mm (IQR 0.45 to 1.38); past it log-odds 2.30, OR 9.93; range across patients 1.86 to 2.73 & 242, 105 & train; test unclear (AUC gain 0.023 [-0.001, 0.055], n=61) & \textcolor{BrickRed}{no} & on the training data only \\
C30 & subgroup: who this part acts on: 105 of 242 patients & 76\% to 91\% (typical patient this part acts on (risk without the part, then with it)) & share 0.43; log-odds 1.03, OR 2.80; range across patients 0.03 to 2.33 & 105, 137 & train; test unclear (AUC gain 0.026 [-0.001, 0.062], n=61) & \textcolor{OliveGreen}{yes} & on the training data only \\
C31 & structure: its formula reproduces it on 242 training patients (largest difference 1.77636e-15) &  & log-odds 2e-15, OR n/a & 242 & train & \textcolor{OliveGreen}{yes} &  \\
C32 & direction: age: 42 years to 66 years (through this part) \newline\textit{shape: hinge at 53 years; 100\% of patients same; verified; one point with C33 (median fold 53 years, typical patient 53 years, 10th-90th percentile 53 years to 53 years)} & 18\% to 19\% & log-odds 0.10, OR 1.11; range across patients 0.10 to 0.10 & 242 & train; test disagrees (AUC gain -0.003 [-0.017, 0.004], n=61) & \textcolor{OliveGreen}{yes} & slightly; on the training data only \\
C33 & threshold: age: fold at about 53 years (starts) \newline\textit{shape: hinge at 53 years; 100\% of patients same; verified; one point with C33 (median fold 53 years, typical patient 53 years, 10th-90th percentile 53 years to 53 years)} & 18\% to 20\% (typical patient with a fold) & fold 52.8 years (IQR 52.8 to 52.8); past it log-odds 0.12, OR 1.13; range across patients 0.12 to 0.12 & 242, 148 & train; test disagrees (AUC gain -0.003 [-0.017, 0.004], n=61) & \textcolor{OliveGreen}{yes} & slightly; on the training data only \\
C34 & subgroup: who this part acts on: 148 of 242 patients & 91\% to 92\% (typical patient this part acts on (risk without the part, then with it)) & share 0.61; log-odds 0.06, OR 1.07; range across patients 0.009 to 0.12 & 148, 94 & train; test disagrees (AUC gain -0.003 [-0.017, 0.004], n=61) & \textcolor{OliveGreen}{yes} & slightly; on the training data only \\
C35 & null: this part barely changes predictions (switched off) &  & max change 0.032 (max $|$change in predicted risk$|$) & 242 & train & \textcolor{BrickRed}{no} &  \\
C36 & structure: its formula reproduces it on 242 training patients (largest difference 0.0) &  & log-odds 0.00, OR n/a & 242 & train & \textcolor{OliveGreen}{yes} &  \\
\bottomrule
\end{tabularx}}
\paragraph{Not supported by the model.}
\begin{itemize}\setlength\itemsep{0pt}\small
\item C17: direction claim about resting ECG in the whole model: not supported by the model (44\% of patients have the opposite sign (allowed: 10\%)).
\item C29: threshold claim about ST depression during exercise (compared with rest) in part ann\_\allowbreak{}27: not supported by the model (the fold changes the curve by only 0.095 log-odds compared with a straight line (needed 0.1)).
\item C35: null claim about age in part ann\_\allowbreak{}28: not supported by the model (predictions move by up to 0.032).
\end{itemize}

\subsection{Coverage}\label{app:heart-s5-c3:coverage}
Coverage is computed by the current code (the function behind the \texttt{get\_coverage} tool) under these conventions. \emph{Candidate nodes} are all nodes of the explained model state except the output nodes. A candidate node is \emph{covered} when some annotation of the hierarchy (a leaf, or a composition together with its descendants) contains it and all its enabled outgoing edges. \emph{Exact Null} parts, which have no effect on any output for any input, count as covered: nodes with no path to an output over enabled, non-zero-weight edges (disconnected inputs and dead ends), and edges that have zero weight or touch such a node. Compositional coverage is the fraction of required compositions present: every annotation except a single global root containing every output must have a parent; a parentless, childless annotation made only of Null nodes is exempt.

\begin{itemize}
\item \textbf{Null, disconnected input (7):} \texttt{-11} (fbs:1), \texttt{-13} (restecg:1), \texttt{-14} (restecg:2), \texttt{-16} (exang:0), \texttt{-20} (slope:2), \texttt{-3} (sex:1), \texttt{-5} (cp:2).
\item \textbf{Uncovered nodes:} none.
\item \textbf{Uncovered edges:} none.
\item \textbf{Missing compositions:} none.
\end{itemize}

\subsection{Annotation hierarchy}\label{app:heart-s5-c3:hierarchy}
The hierarchy has 1 root; the global root is heart\_\allowbreak{}disease\_\allowbreak{}logit. Tags (A1, A2, \ldots) are used in the figures.
\begingroup\parindent=0pt\parskip=1pt
\hspace*{0.0em}$\triangleright$~A1: heart\_\allowbreak{}disease\_\allowbreak{}logit (composition, \cref{app:heart-s5-c3:ann-heart-disease-logit})\par
\hspace*{1.5em}$\triangleright$~A2: ischaemia\_\allowbreak{}chain (leaf, \cref{app:heart-s5-c3:ann-ischaemia-chain})\par
\hspace*{1.5em}$\triangleright$~A3: age\_\allowbreak{}hinge (leaf, \cref{app:heart-s5-c3:ann-age-hinge})\par
\endgroup

\subsection{Network figures}\label{app:heart-s5-c3:figures}
\begin{figure}[htbp]
\centering
\input{\explaneatcertdir{heart-s5-c3}heart-s5-c3-fig-full.tex}
\caption{heart-s5-c3: full view of the explained model state.}
\label{app:heart-s5-c3:fig-full}
\end{figure}
\begin{figure}[htbp]
\centering
\input{\explaneatcertdir{heart-s5-c3}heart-s5-c3-fig-level-0.tex}
\caption{heart-s5-c3: collapsed view at hierarchy level 0 (fully recomposed).}
\label{app:heart-s5-c3:fig-level-0}
\end{figure}
\begin{figure}[htbp]
\centering
\input{\explaneatcertdir{heart-s5-c3}heart-s5-c3-fig-level-1.tex}
\caption{heart-s5-c3: collapsed view at hierarchy level 1.}
\label{app:heart-s5-c3:fig-level-1}
\end{figure}

\subsection{Refactoring operations}\label{app:heart-s5-c3:operations}
The explained model state is the phenotype after these operations, replayed in order. All are function-preserving; the largest change in the model output over the 242 training rows is $2.22\times 10^{-16}$ (tolerance $1.00\times 10^{-9}$).
\begin{itemize}
\item \#26: split \texttt{-12} into \texttt{-12\_\allowbreak{}a}, \texttt{-12\_\allowbreak{}b}.
\end{itemize}

\subsection{Annotations}\label{app:heart-s5-c3:annotations}
\subsubsection{\texorpdfstring{A1: heart\_\allowbreak{}disease\_\allowbreak{}logit}{A1: heart disease logit}}\label{app:heart-s5-c3:ann-heart-disease-logit}
% annotation id ann_29, operation seq 29
\textbf{Kind:} Composition annotation (depth 0); children A2 (\cref{app:heart-s5-c3:ann-ischaemia-chain}), A3 (\cref{app:heart-s5-c3:ann-age-hinge}).

\begin{itemize}
\item \textbf{Entry nodes:} \texttt{-26} (thal:7), \texttt{-25} (thal:6), \texttt{-24} (thal:4), \texttt{-23} (thal:3), \texttt{-22} (ca), \texttt{-21} (slope:3), \texttt{-20} (slope:2), \texttt{-19} (slope:1), \texttt{-17} (exang:1), \texttt{-16} (exang:0), \texttt{-15} (thalach), \texttt{-14} (restecg:2), \texttt{-13} (restecg:1), \texttt{-11} (fbs:1), \texttt{-10} (fbs:0), \texttt{-9} (chol), \texttt{-8} (trestbps), \texttt{-7} (cp:4), \texttt{-6} (cp:3), \texttt{-5} (cp:2), \texttt{-4} (cp:1), \texttt{-3} (sex:1), \texttt{-2} (sex:0), \texttt{-12\_\allowbreak{}a} (restecg:0\_\allowbreak{}a).
\item \textbf{Exit nodes:} \texttt{0}.
\item \textbf{Internal nodes:} none.
\item \textbf{Edges (weight):} \texttt{-2}$\to$\texttt{0}~(-1.04), \texttt{-4}$\to$\texttt{0}~(-0.494), \texttt{-6}$\to$\texttt{0}~(-0.442), \texttt{-7}$\to$\texttt{0}~(0.354), \texttt{-8}$\to$\texttt{0}~(0.483), \texttt{-9}$\to$\texttt{0}~(0.551), \texttt{-10}$\to$\texttt{0}~(0.274), \texttt{-15}$\to$\texttt{0}~(-0.213), \texttt{-17}$\to$\texttt{0}~(0.423), \texttt{-19}$\to$\texttt{0}~(-0.873), \texttt{-21}$\to$\texttt{0}~(-0.434), \texttt{-22}$\to$\texttt{0}~(1.49), \texttt{-23}$\to$\texttt{0}~(-0.471), \texttt{-24}$\to$\texttt{0}~(0.172), \texttt{-25}$\to$\texttt{0}~(-0.421), \texttt{-26}$\to$\texttt{0}~(0.0253), \texttt{204}$\to$\texttt{0}~(0.0487), \texttt{290}$\to$\texttt{0}~(0.216), \texttt{1987}$\to$\texttt{0}~(0.0824), \texttt{-12\_\allowbreak{}a}$\to$\texttt{0}~(-0.188).
\item \textbf{With its descendants:} 32 nodes.
\item \textbf{Covers} (coverage, own and as part of the hierarchy): 31 candidate nodes.
\end{itemize}

\textbf{Hypothesis.} Root: p = sigmoid(z), z = -0.536 + LINEAR CORE + 0.049*\allowbreak{}ischaemia\_\allowbreak{}chain[204]\allowbreak{} + 0.216*\allowbreak{}ischaemia\_\allowbreak{}chain[290]\allowbreak{} + 0.082*age\_\allowbreak{}hinge. LINEAR CORE (17 direct input-$>$output edges, z-scored inputs): ca +1.485, sex:0 -1.042, slope:1 -0.873, chol +0.551, cp:1 -0.494, trestbps +0.483, thal:3 -0.471, cp:3 -0.442, slope:3 -0.434, exang:1 +0.423, thal:6 -0.421, cp:4 +0.354, fbs:0 +0.274, thalach -0.213, restecg:0 -0.188, thal:4 +0.172, thal:7 +0.025. Seven inputs are unconnected (sex:1, cp:2, fbs:1, restecg:1, restecg:2, exang:0, slope:2; zero-effect entries of this root). Oldpeak and age have no direct edge and act only through the children; restecg:0 acts both directly and through the ischaemia chain (split into restecg:0\_\allowbreak{}a direct and restecg:0\_\allowbreak{}b). COMPOSITION: additive. The two children are read out additively into z: ischaemia\_\allowbreak{}chain (two exits: a one-sided ST-depression score opened earlier by a normal resting ECG, and a steeper second stage of the same score) and age\_\allowbreak{}hinge (an age term that starts just below the median age). Together the chain gives oldpeak a hinge-shaped effect: about 0.12 logit per SD between its two thresholds and about 1.10 per SD above the second (z-oldpeak -0.39 with a normal ECG, +0.40 otherwise). EVIDENCE (train, n=242): full model AUC 0.9376, acc 0.888 (compute\_\allowbreak{}performance)\allowbreak{}. Linear core alone AUC 0.9317; logit SD 3.22 (core) vs 0.80 (nonlinear readouts). Removing the whole chain (204 with downstream) gives 0.9317; readout-only ablations: 204 0.9380, 290 0.9341, age 0.9371. All ReLUs as identities 0.9304; one least-squares linear fit of all nonlinear readouts on restecg:0, oldpeak and age gives 0.9355 (R2 0.88); an L2 logistic regression on all 26 inputs gives 0.9355. Mean $|$SHAP$|$ (kernel, all 242 train rows): ca 0.141, sex:0 0.104, slope:1 0.099, oldpeak 0.060, thal:3 0.054, exang:1 0.045, chol 0.044, cp:3 0.043, cp:4 0.040, trestbps 0.040. Interpretation: close to logistic regression plus one ST-depression correction that switches on above a threshold (lower for patients with a normal resting ECG) and a small age term above median age. Counter-directional signs reported, not explained: a normal resting ECG lowers the oldpeak threshold inside the chain (raising risk), fbs:0 (normal fasting sugar) +0.27 directly, thal:7 (reversible defect) only +0.025. Compared with the 10-seed heart-s5 record (generation 37, key 9330: ischaemia gate -$>$ resting-BP gate -$>$ age term) and its cluster-1 sibling: the same ischaemia gate on oldpeak and restecg:0, but here the second stage is a steeper copy of the score (no resting-BP offset) and the age term is a free-standing hinge (not switched off by the chain).

\textbf{Closed form} (children as named functions):
{\small\begingroup\allowdisplaybreaks
\begin{align*}
y_{0} &= \sigma\bigl(z_{0}\bigr) \\
z_{0} &= -0.536 + 0.216 \operatorname{ischaemia\_chain\_1}\bigl(x_{restecg:0\_b},  \\
&\qquad x_{oldpeak}\bigr) + 0.354 x_{cp:4} + 0.551 x_{chol} + 0.172 x_{thal:4} + 1.49 x_{ca} \\
&\qquad + 0.0487 \operatorname{ischaemia\_chain\_0}\bigl(x_{restecg:0\_b}, x_{oldpeak}\bigr) \\
&\qquad + 0.274 x_{fbs:0} + 0.0253 x_{thal:7} + 0.483 x_{trestbps} \\
&\qquad + 0.0824 \operatorname{age\_hinge}\bigl(x_{age}\bigr) + 0.423 x_{exang:1} - 0.213 x_{thalach} \\
&\qquad - 0.494 x_{cp:1} - 1.04 x_{sex:0} - 0.188 x_{restecg:0\_a} - 0.421 x_{thal:6} \\
&\qquad - 0.471 x_{thal:3} - 0.434 x_{slope:3} - 0.442 x_{cp:3} - 0.873 x_{slope:1}
\end{align*}\endgroup}

Expanded over the children:
{\small\begingroup\allowdisplaybreaks
\begin{align*}
y_{0} &= \sigma\bigl(z_{0}\bigr) \\
z_{0} &= -0.536 + 0.216 \operatorname{relu}\bigl(-0.468 \\
&\qquad + 0.862 \operatorname{relu}\bigl(-0.00219 + 0.94 \operatorname{relu}\bigl(0.245 \\
&\qquad + 2.5 x_{oldpeak} + 0.98 x_{restecg:0\_b}\bigr)\bigr) + 1.01 \operatorname{relu}\bigl(0.245 \\
&\qquad + 2.5 x_{oldpeak} + 0.98 x_{restecg:0\_b}\bigr)\bigr) + 0.354 x_{cp:4} + 0.551 x_{chol} \\
&\qquad + 0.172 x_{thal:4} + 1.49 x_{ca} + 0.0487 \operatorname{relu}\bigl(0.245 + 2.5 x_{oldpeak} \\
&\qquad + 0.98 x_{restecg:0\_b}\bigr) + 0.274 x_{fbs:0} + 0.0253 x_{thal:7} + 0.483 x_{trestbps} \\
&\qquad + 0.0824 \operatorname{relu}\bigl(0.191 + 0.872 x_{age}\bigr) + 0.423 x_{exang:1} \\
&\qquad - 0.213 x_{thalach} - 0.494 x_{cp:1} - 1.04 x_{sex:0} - 0.188 x_{restecg:0\_a} \\
&\qquad - 0.421 x_{thal:6} - 0.471 x_{thal:3} - 0.434 x_{slope:3} - 0.442 x_{cp:3} - 0.873 x_{slope:1}
\end{align*}\endgroup}

where $y_{0}$ is node \texttt{0}; $z_{i}$, where shown, is the pre-activation of a sigmoid output $y_{i}$; weights and biases are shown to 3 significant figures.
 On the training split (n = 242), the displayed closed form differs from the forward pass by at most 0.00216, and the exact function by at most $2.78\times 10^{-16}$.

\textbf{Evidence} (2 items; full payloads in the JSON record):
\begin{itemize}
\item \textsc{analytical}, \texttt{analytical\_\allowbreak{}formula}; data split: n/a (analytical). Closed form of the full model and its collapsed form in terms of the children.
\item \textsc{empirical}, \texttt{empirical\_\allowbreak{}viz}; data split: train. Train-split performance, whole-model SHAP and ablation / linear-replacement analysis.
\end{itemize}

\textbf{Authorship:} drafted by AI-assisted draft on 2026-11-10; verified by nobody (unverified).

\subsubsection{\texorpdfstring{A2: ischaemia\_\allowbreak{}chain}{A2: ischaemia chain}}\label{app:heart-s5-c3:ann-ischaemia-chain}
% annotation id ann_27, operation seq 27
\textbf{Kind:} Leaf annotation (depth 1); parent A1 (\cref{app:heart-s5-c3:ann-heart-disease-logit}).

\begin{itemize}
\item \textbf{Entry nodes:} \texttt{-12\_\allowbreak{}b} (restecg:0\_\allowbreak{}b), \texttt{-18} (oldpeak).
\item \textbf{Exit nodes:} \texttt{204}, \texttt{290}.
\item \textbf{Internal nodes:} \texttt{3487}.
\item \textbf{Edges (weight):} \texttt{-12\_\allowbreak{}b}$\to$\texttt{204}~(0.98), \texttt{-18}$\to$\texttt{204}~(2.5), \texttt{204}$\to$\texttt{3487}~(0.94), \texttt{204}$\to$\texttt{290}~(1.01), \texttt{3487}$\to$\texttt{290}~(0.862).
\item \textbf{Covers} (coverage, own and as part of the hierarchy): 3 candidate nodes.
\end{itemize}

\textbf{Hypothesis.} Ischaemia chain (raises risk with ST depression, in two stages). 204 = ReLU(0.245 + 2.501*oldpeak + 0.980*restecg:0); 3487 = ReLU(0.940*204 - 0.002) (a near-copy of 204); 290 = ReLU(-0.468 + 1.008*204 + 0.862*3487) = ReLU(1.818*204 - 0.470) whenever 204 $>$ 0.0023. Readouts: logit += 0.049*204 + 0.216*290. Hypothesis: an ST-depression score that opens above an oldpeak threshold set by the resting ECG (z-oldpeak -0.49 with a normal ECG, +0.29 otherwise) and then a second, much steeper stage once the score passes 0.259 (z-oldpeak -0.39 with a normal ECG, +0.40 otherwise). Net effect on the logit: about 0.12 per SD of oldpeak between the two thresholds and about 1.10 per SD above the second, i.e. a hinge in oldpeak whose position moves about 0.78 SD lower with a normal resting ECG. EVIDENCE (train, n=242): 204 active for 105/242 (58 of 121 with a normal ECG, 47 of 121 without); 290 active for 96/242; disease rate 66\% when 290 is active vs 30\% when not; corr(204, oldpeak) 0.91. Sensitivities (mean $|$d exit/d entry$|$): 290: oldpeak 1.82, restecg:0 0.71; 204: oldpeak 1.09, restecg:0 0.43. Ablation: removing 204 with its downstream (the whole chain) lowers train AUC 0.9376 -$>$ 0.9317, the linear-core level; removing the steepener readout (290) alone 0.9341; removing 204's own readout 0.9380 (no loss). An affine replacement of 290 on 204 fits almost perfectly (R2 0.998, AUC 0.9376): the steepening is a hinge that the data barely probe on its flat side. Supported. Oldpeak has no direct edge and acts only here. Counter-intuitive sign reported, not explained: a normal resting ECG (restecg:0) lowers the oldpeak threshold (raising risk) inside the chain, partly offset by its direct -0.188 edge.

\textbf{Closed form:}
{\small\begingroup\allowdisplaybreaks
\begin{align*}
y_{0} &= \operatorname{relu}\bigl(0.245 + 2.5 x_{oldpeak} + 0.98 x_{restecg:0\_b}\bigr) \\
y_{1} &= \operatorname{relu}\bigl(-0.468 + 0.862 \operatorname{relu}\bigl(-0.00219 \\
&\qquad + 0.94 \operatorname{relu}\bigl(0.245 + 2.5 x_{oldpeak} + 0.98 x_{restecg:0\_b}\bigr)\bigr) \\
&\qquad + 1.01 \operatorname{relu}\bigl(0.245 + 2.5 x_{oldpeak} + 0.98 x_{restecg:0\_b}\bigr)\bigr)
\end{align*}\endgroup}

where $y_{0}$ is node \texttt{204}; $y_{1}$ is node \texttt{290}; $z_{i}$, where shown, is the pre-activation of a sigmoid output $y_{i}$; weights and biases are shown to 3 significant figures.
 On the training split (n = 242), the displayed closed form differs from the forward pass by at most 0.0205, and the exact function by at most $1.78\times 10^{-15}$.

\textbf{Evidence} (2 items; full payloads in the JSON record):
\begin{itemize}
\item \textsc{analytical}, \texttt{analytical\_\allowbreak{}formula}; data split: n/a (analytical). Closed form from get\_\allowbreak{}formula.
\item \textsc{empirical}, \texttt{empirical\_\allowbreak{}viz}; data split: train. Train-split sensitivity, line plot, activity, outcome rates and ablations.
\end{itemize}

\textbf{Authorship:} drafted by AI-assisted draft on 2026-11-10; verified by nobody (unverified).

\subsubsection{\texorpdfstring{A3: age\_\allowbreak{}hinge}{A3: age hinge}}\label{app:heart-s5-c3:ann-age-hinge}
% annotation id ann_28, operation seq 28
\textbf{Kind:} Leaf annotation (depth 1); parent A1 (\cref{app:heart-s5-c3:ann-heart-disease-logit}).

\begin{itemize}
\item \textbf{Entry nodes:} \texttt{-1} (age).
\item \textbf{Exit nodes:} \texttt{1987}.
\item \textbf{Internal nodes:} none.
\item \textbf{Edges (weight):} \texttt{-1}$\to$\texttt{1987}~(0.872).
\item \textbf{Covers} (coverage, own and as part of the hierarchy): 1 candidate node.
\end{itemize}

\textbf{Hypothesis.} Age hinge (raises risk with age above a threshold, small). 1987 = ReLU(0.191 + 0.872*age), read out +0.082. Hypothesis: age's only route to the output; nothing below z-age -0.22 (just under the median age), then +0.072 logit per SD of age. EVIDENCE (train, n=242): active for 148/242; disease rate 55\% when active vs 28\% when not; line plot correlation 0.90; mean $|$sensitivity$|$ 0.53. Removing it lowers train AUC 0.9376 -$>$ 0.9371; an affine replacement (R2 0.80) gives 0.9376. Supported, near-negligible.

\textbf{Closed form:}
{\small\begingroup\allowdisplaybreaks
\begin{align*}
y_{0} &= \operatorname{relu}\bigl(0.191 + 0.872 x_{age}\bigr)
\end{align*}\endgroup}

where $y_{0}$ is node \texttt{1987}; $z_{i}$, where shown, is the pre-activation of a sigmoid output $y_{i}$; weights and biases are shown to 3 significant figures.
 On the training split (n = 242), the displayed closed form differs from the forward pass by at most $2.75\times 10^{-4}$, and the exact function by at most 0.

\textbf{Evidence} (2 items; full payloads in the JSON record):
\begin{itemize}
\item \textsc{analytical}, \texttt{analytical\_\allowbreak{}formula}; data split: n/a (analytical). Closed form from get\_\allowbreak{}formula.
\item \textsc{empirical}, \texttt{empirical\_\allowbreak{}viz}; data split: train. Train-split sensitivity, line plot, activity, outcome rates and ablations.
\end{itemize}

\textbf{Authorship:} drafted by AI-assisted draft on 2026-11-10; verified by nobody (unverified).

\subsection{Audit against the definitions}\label{app:heart-s5-c3:audit}
Automatic findings on the stored explanation (codes as in the JSON record's \texttt{issues}); they flag what an examiner should check, not verdicts.
\begin{itemize}
\item \textsc{warning} \texttt{unverified}: No human verification is recorded for this explanation.
\end{itemize}
\end{drafted}
